\documentclass[11pt,a4paper]{article}
\usepackage{fontspec}
\setmainfont{Noto Serif}
\setsansfont{Noto Sans}
\usepackage{babel}
\babelprovide[main,import]{russian}
\usepackage{amsmath,amssymb,mathtools,geometry,fancyhdr,microtype,hyperref,xcolor}
\geometry{left=21mm,right=21mm,top=20mm,bottom=20mm}
\definecolor{ink}{HTML}{173D57}
\hypersetup{colorlinks=true,urlcolor=ink,pdftitle={BGS-F. Интерференция при нескольких совмещениях}}
\pagestyle{fancy}\fancyhf{}\lhead{\small BGS-F}\rhead{\small Несколько совмещений траекторий}\cfoot{\thepage}
\setlength{\headheight}{15pt}\setlength{\parskip}{4pt}
\setlength{\emergencystretch}{3em}\tolerance=3000
\newcommand{\Tr}{\operatorname{Tr}}
\newcommand{\C}{\mathbb C}
\newcommand{\Z}{\mathbb Z}
\newcommand{\sect}[1]{\section*{\sffamily\color{ink}#1}}
\newcommand{\subt}[1]{\par\medskip\noindent\textbf{#1.}\quad}
\begin{document}
\raggedright
\noindent{\small\sffamily\color{ink} МАТЕМАТИЧЕСКАЯ ФИЗИКА\hfill 22 НОЯБРЯ 2026}
\begin{center}{\LARGE\sffamily Несколько совмещений траекторий}\par\vspace{4pt}
{\large Снятие повторного счёта до порога строгого большинства}
\end{center}

Для фиксированной карты \(U_N=e^{i10^{-80}NA_N}M_N\) обозначим
\(a(x)=\prod_{t=0}^{m-1}(U_N)_{x_{t+1},x_t}\), \(x_m=x_0\).
С каждой парой \(x,y\) сравниваем \(x_t\) с \(y_{s+t}\) и \(y_{s-t}\), при всех \(s\bmod m\).
Пусть \(H_{m,N}(d)\) есть сумма \(a(x)\overline{a(y)}\) по парам, для которых хотя бы одно сравнение даёт не менее \(dm\) совпадений. Пара входит в сумму один раз.

\subt{Доказанная теорема}
Для каждого фиксированного \(1/2<d\le1\) существуют \(c>0\) и \(N_0\), такие что при всех нечётных \(N\ge N_0\) и \(m\ge\lceil\log N\rceil\)
\[
           H_{m,N}(d)=m+O(e^{-cm}).
\]
После удаления циклических копий \(y\) слова \(x\) остаётся \(O(e^{-cm})\) (с возможным уменьшением \(c\)). Это расширяет прежний диапазон \(d>2/3\) на все \(d>1/2\), одновременно для прямых и обратных совмещений. Доказательство относится к исходным комплексным амплитудам и охватывает чётные и составные \(m\).

На границе \(d=1/2\) для двух прямых совмещений, сдвинутых на \(m/2\), отдельно вычислен размер вклада без тройных совпадений. Он равен \(O(N^{-1})\) на одно такое пересечение при \(m\asymp N\). Доказательство опирается на явную матрицу переключения двух парных схем.

\subt{Как устраняются повторные отметки}
Если два совмещения одного направления разнесены сдвигом большого порядка, работает счёт числа свободных букв. Остальные пересечения складываются в петлю с ограниченным числом параллельных нитей. Превышение половины всех возможных совпадений вынуждает тройное совпадение на положительной доле сечений. Для поглощающего сжатия оператор возврата имеет степень, стремящуюся к нулю. Это позволяет оценить полную знаковую сумму пересечения и провести включение--исключение по всем отметкам.

\subt{Что теперь требуется для парного предела}
Пусть \(R_{m,N}(d)\) есть знаковая сумма по парам с числом совпадений \(<dm\) при \emph{каждом} прямом и обратном совмещении. Для исходного веса \(w\) доказано
\[
 \Delta_{2,N}(b_*)=\delta_*+
     \frac1{N^2}\sum_{\lceil\log N\rceil\le m<2N} w(m/N)R_{m,N}(d)+o(1),
      \qquad \delta_*=\int_1^2(t-1)w(t)\,dt .
\]
Таким образом, вся отрицательная поправка должна поступить из класса с меньшим числом совпадений. Предел этого остатка и локальный закон \(\mathrm{Sine}_2\) остаются открытыми.

\newpage
\sect{1. Входные оценки и один фиксированный сдвиг}
В доказательстве используются три уже установленных для исходного \(U_N\) факта из предыдущих глав сводного отчёта.

\begin{enumerate}
\item Для каждого фиксированного \(j\ge1\),
\(\max_{a,b}|(U_N^j)_{ab}|\le C_j N^{-1/2}\).
Для сопряжения, транспонирования и обратной степени верна та же оценка. Для \(j=1\) обозначим через \(K\) верхнюю постоянную в \(\max |(U_N)_{ab}|^2\le K/N\).
\item Для \(V_+=U_N\otimes\overline{U_N}\) и \(V_-=U_N\otimes U_N^*\) обозначим через \(P\) проектор на \(x=y\). В прямом случае ограничим пространство на \(e^\perp\), \(e=N^{-1/2}\sum_x|x,x\rangle\), и используем \(P_0=P-ee^*\). В обратном случае \(P_0=P\). Тогда
\[
 \sum_{j\ge1}\limsup_{N\to\infty}
    \|P_0 V_+^jP_0\|\le10^{-5},\qquad
 \lim_{N\to\infty}\|P V_-^jP\|=0\quad\text{при каждом }j.
\]
В частности, \(\||U_N|^{\circ2}-J_N\|<1/100\) при больших \(N\), \(J_N(x,y)=1/N\).
\item \(\sum_{1\le m<\lceil\log N\rceil}|\Tr U_N^m|^2/N^2=o(1)\).
Доказанная через давление полоса в действительности длиннее \(2\log N/\log(3+2\sqrt2)\).
\end{enumerate}

Первые два факта взяты из главы о контактах, формулы (4), (5), (7), с разбором стационарных ветвей при каждом фиксированном числе шагов. Третий факт взят из главы о давлении, формула (1). Никакой входной сходимости парной корреляции или предположения о случайности фаз не требуется.

Для одной прямой отметки производящая функция числа совпадений есть
\(\Tr([(I+(z-1)P)V_+]^m)\).
Для одной обратной отметки заменяем \(V_+\) на \(V_-\): обращение порядка второго слова даёт транспонирование его покомпонентно сопряжённой матрицы.
В прямом случае линия \(e\) даёт член \(z^m\) точно. На остальных пространствах применяем резольвенту \(V_\pm\).
При \(|z|=100\) и \(|\lambda|\ge e^{1/20}\) возможный полюс потребовал бы необратимости
\[
        I-(z-1)\sum_{j\ge1}\lambda^{-j}P_0V_\pm^j P_0 .
\]
Норма поправки меньше 1 при больших \(N\). Хвост ряда равномерен по \(N\) по унитарности, а первые члены ограничены вторым входным фактом. След степени на остаточном пространстве по собственным числам ограничен \(N^2 e^{m/20}\); диагонализуемость для следа не нужна.
По Коши сумма коэффициентов при \(h\ge k=\lceil dm\rceil\) отличается от \(1\) для прямого и от \(0\) для обратного сравнения не более чем на
\[
              \frac{N^2}{1-1/100}\,e^{m/20}100^{-k}.
\]
После суммирования по \(2m\) сравнениям при \(m\ge\lceil\log N\rceil\) это \(O(e^{-c_0 m})\):
\(d\log100-1/20-2>0\) даже при \(d=1/2\).

\newpage
\sect{2. Далёкие пересечения и конечное число нитей}
Сравнение обозначаем \(f(t)=t+s\) или \(f(t)=-t+s\) на \(\Z/m\Z\).
В событии \(E_f\) хотя бы \(k=\lceil dm\rceil\) равенств \(x_t=y_{f(t)}\).
Для множества \(I\) сравнений пишем \(E_I=\bigcap_{f\in I} E_f\).

Если \(f_1,f_2\) одного направления, то \(f_1^{-1}f_2\) есть сдвиг. Пусть его порядок \(r\ge2\). Полное объединение двух наборов равенств состоит из циклов длины \(2r\) в двудольном графе на позициях слов \(x,y\). При фиксированном выборе \(k\) выполненных равенств из каждого набора число целых циклов не превосходит \(k/r\). Ранг отождествлений равен числу рёбер минус число циклов. Поэтому для данного пересечения
\[
 \sum_{(x,y)\in E_{\{f_1,f_2\}}}|a(x)\overline{a(y)}|
       \le \binom{m}{k}^2 K^m N^{m-2k+k/r}.                  \tag{1}
\]
Здесь каждый из \(2m\) квантовых множителей имеет модуль не больше
\(\sqrt{K/N}\), а свободная буквенная метка даёт один множитель \(N\).

Выберем фиксированное \(r_0\ge2\) с
\(2d-1-d/(r_0+1)>0\), и положим \(L_0=\operatorname{lcm}(1,\ldots,r_0)\).
Назовём пересечение большим по порядку, если в нём есть два сравнения одного направления с относительным вращением порядка \(>r_0\). Для каждого такого множества \(I\) сумма по \(E_I\) по модулю не превосходит правой части (1) для этой пары. Даже после умножения на \(2^{2m}\), число всех подмножеств сравнений, оценка даёт \(O(e^{-Am})\) для любого фиксированного \(A>0\) при больших \(N\).

Рассмотрим оставшиеся множества \(I\). Положим \(r=\gcd(L_0,m)\), \(b=m/r\). Относительные вращения внутри каждого направления лежат в подгруппе сдвигов кратных \(b\): их порядки делят \(L_0\). Выберем \(f_0\in I\), предпочтительно прямое, если оно есть, и заменим \(y_t\) на \(y_{f_0(t)}\). Теперь все сравнения в новой координате имеют один из видов
\[
                    t\longmapsto t+ib,\qquad
                    t\longmapsto c-t+ib .
\]
Второй вид появляется только при наличии двух направлений, а \(c\) один для всего \(I\) по модулю \(b\). В частности, \(|I|\le2r\le2L_0\).
Число таких подмножеств ограничено \((2m+1)^{2L_0+1}\).

Для каждого \(I\) с \(|I|\ge2\) фиксируем два сравнения из него. Далее сохраняем одновременно \emph{все} метки из \(I\). Оценка только знаковой суммы двойных пересечений сама по себе не допускает снятия кратности. Граница по порядку нужна для контроля каждого члена полного включения--исключения.

\newpage
\sect{3. Сложение временной окружности}
После выбора \(I\) все его сравнения становятся одновременными равенствами между нитями конечной ширины. Вес остаётся тем же произведением элементов.

\subt{Только сдвиги}
Сечение с номером \(j\), \(0\le j<b\), содержит позиции \(j+ib\), \(0\le i<r\), в каждом из двух слов. Соседние сечения соединяет \(W\), тензорное произведение \(2r\) матриц, каждая из которых \(U_N\), \(\overline{U_N}\) или их транспонирование. При переходе через конец окружности меняется номер нити. Эту перестановку учитываем в граничной матрице. Число позиций одного слова в одном сечении равно \(v=r\).

\subt{Сдвиги и отражения}
Берём целые представители \(j\) внутри
\[
                         c/2< j <(c+b)/2 .
\]
Сечение содержит позиции \(j+ib\) и \(c-j+ib\), \(0\le i<r\), в обоих словах. Они различны. При изменении \(j\) одни нити идут вперёд, другие назад. Матрица \(W\) теперь содержит \(4r\) тензорных множителей, включая транспонированные для обратных нитей. Строгий интервал выбирает один представитель каждой двухэлементной орбиты \(j\mapsto c-j\bmod b\). Вне сечений остаются только неподвижные классы этой инволюции, не более двух, то есть не более \(2r\) позиций каждого слова. В каждом сечении \(v=2r\). Суммирование по оставшимся позициям и умножение на граничные амплитуды дают одну матрицу между концами сложенного отрезка.

В обоих случаях число регулярных сечений \(n=m/v+O_{L_0}(1)\).
Каждому \(f\in I\) поставим в соответствие переменную \(z_f\), отмечающую выполненные равенства. Пусть \(D(\mathbf z)\) обозначает диагональную матрицу произведения этих отметок на одном регулярном сечении. Она унитарна при \(|z_f|=1\). Каждая отметка имеет вид \(I+(z_f-1)P_e\), где \(P_e\) есть диагональный проектор на равенство букв двух нитей \(e\).

Два выбранных сравнения задают в регулярном сечении два совершенных паросочетания без общего ребра. Для одинакового направления совпадение рёбер означало бы нулевой относительный сдвиг. Для разных направлений оно потребовало бы фиксированной позиции относительного отражения, а такие позиции вынесены из регулярных сечений. Если выполнено больше \(v\) из этих \(2v\) равенств, некоторые два ребра имеют общий конец. Три разные нити тогда несут одну букву. Пусть \(R\) есть проектор на событие «выполнено больше \(v\) равенств». Вводим ещё одну переменную \(u\) для числа таких сечений, кроме первого регулярного сечения.

Точный многочлен всех отметок записывается
\[
       {\cal F}_I(\mathbf z,u)
           =\Tr\!\left(B(\mathbf z)
             [\{I+(u-1)R\}D(\mathbf z) W]^{n-1}\right).        \tag{2}
\]
Для проверки разложим след по состояниям \(v_0,\ldots,v_{n-1}\). Фактор
\(D(\mathbf z)\) действует на выходе каждого шага. В \(B\) входят отметки первого сечения, условия замыкания, дополнительные амплитуды у краёв и сумма по нерегулярным позициям. Все эти факторы зависят лишь от двух крайних сечений и нерегулярных букв. При \(|z_f|=1\) каждая амплитуда по модулю не больше 1, поэтому
\(\|B(\mathbf z)\|_1\le N^{12r}\).
Здесь \(\|\cdot\|_1\) есть следовая норма: размер не больше \(N^{4r}\), а каждый элемент — сумма не более \(N^{4r}\) слагаемых. Формула не требует коммутирования \(B\) с переносом.

\newpage
\sect{4. Норма после встречи с тройным совпадением}
В этом разделе размер тензора \(p\) фиксирован, \(W=\bigotimes_{i=1}^p W_i\), каждый \(W_i\) является одним из вариантов \(U_N\).
Пусть \(P_A\) проектирует на совпадение букв всех нитей из тройки \(A\).
Рассмотрим цепочку из \(j\ge1\) шагов \(W\) после \(P_A\). Пусть до последнего шага стоят только равенства между нитями вне \(A\), а в конце стоит равенство \(P_e\), где двухэлементное \(e\) касается \(A\). Тогда норма всей цепочки ограничена \(C_j N^{-1/2}\).

Действительно, вход имеет вид \(\sum_x |x\rangle^{\otimes A}\otimes \xi_x\),
\(\sum_x\|\xi_x\|^2=1\). До последней отметки перенос факторизуется в
\(\bigotimes_{i\in A} W_i^j\) и одно сжатие \(G\) на остальных нитях,
\(\|G\|\le1\). Выберем \(k\in A\setminus e\). Векторы \(W_k^j|x\rangle\) ортогональны для разных \(x\), и эта нить свободна в последней отметке. Значит, нормы после неё складываются по \(x\) квадратично.
Если оба конца \(e\) лежат в \(A\), для каждого \(x\) возникает множитель
\[
       \sum_y |(W_{e_1}^j)_{y x}|^2 |(W_{e_2}^j)_{y x}|^2
                      \le C_j^2/N.
\]
Если в \(A\) один конец, для каждого \(x\) сумма квадратов равна
\(\sum_y |(W_{e_1}^j)_{y x}|^2\|\Pi_{e_2,y} G\xi_x\|^2\), где \(\Pi_{e_2,y}\) фиксирует букву \(y\) на внешней нити. Она не больше
\(C_j^2\|\xi_x\|^2/N\). Это доказывает оценку.

Проектор \(R\) из (2) раскладывается конечным включением--исключением:
\[
 R=\sum_{\substack{S\subset E\\|S|\ge v+1}}
     (-1)^{|S|-v-1}\binom{|S|-1}{v}\prod_{e\in S}P_e,       \tag{3}
\]
где \(E\) — два выбранных паросочетания. Это разложение индикатора порога на булевом кубе. В каждом ненулевом слагаемом уже имеется проектор на тройное совпадение \(P_A\), по свойству \(v+1\) рёбер.

Пишем \(V=D(\mathbf z)W\). Для любого фиксированного набора положительных \(j_1,\ldots,j_h\) при \(h=p\) утверждаем
\[
             \|R V^{j_h} R\cdots R V^{j_1} R\|=o(1)
             \quad (|z_f|=1) .                               \tag{4}
\]
При раскрытии каждого \(R\) по (3), а каждого \(D\) по его конечному произведению, получаем ограниченное по числу слагаемых разложение, с коэффициентами, ограниченными независимо от \(N,\mathbf z\).
В каждом слагаемом среди \(h+1\) тройных совпадений, выбранных в разных временах, какие-то два касаются общей нити. Рассмотрим более раннюю из этих троек и первое последующее равенство, касающееся её нитей. Такое равенство обязательно существует. Другие равенства в том же сечении коммутируют с ним. Между этой тройкой и выбранным равенством все равенства лежат вне её нитей. Поэтому доказанная оценка применяется с \(1\le j\le\sum_\ell j_\ell\).
Остальные факторы имеют норму не больше 1. Это даёт в (4) ошибку \(O(N^{-1/2})\) для фиксированных \(j_1,\ldots,j_h\).

\newpage
\sect{5. Резольвента и экспоненциальная оценка пересечений}
Для \(|\lambda|\ge2\) и \(|z_f|=1\) введём на \(\operatorname{ran}R\) сжатый ряд
\[
      {\cal L}(\lambda)
        =R V(\lambda-V)^{-1} R
        =\sum_{j\ge1}\lambda^{-j} R V^j R .
\]
Его норма не превосходит 1. Утверждение (4) влечёт
\(\sup\|{\cal L}(\lambda)^h\|=o(1)\). Для этого сперва обрезаем каждый ряд на произвольном \(J\). Разность степеней не больше \(h\,2^{-J}\) по телескопированию, а каждый член конечного разложения стремится к нулю равномерно. Затем увеличиваем \(J\). Утверждение равномерно для всех типов сложения ограниченной ширины, так как число возможных сечений и наборов паросочетаний ограничено в зависимости только от \(L_0\).

Пусть \(A(u)=[I+(u-1)R]V\) и \(|u|\le M\), \(M>1\) фиксировано.
Для достаточно больших \(N\) норма \([(u-1){\cal L}(\lambda)]^h\) меньше \(1/2\). Конечная геометрическая сумма с последующим обращением степени даёт обратимость и равномерную оценку нормы
\([I-(u-1){\cal L}(\lambda)]^{-1}\).
Факторизация
\[
 \lambda-A(u)=
       [I-(u-1)R V(\lambda-V)^{-1}](\lambda-V)
\]
переносит эту оценку на \((\lambda-A(u))^{-1}\). В самом деле, у оператора
\(R V(\lambda-V)^{-1}\) образ лежит в \(\operatorname{ran}R\), а его сжатие на этом пространстве есть \({\cal L}(\lambda)\); относительно \(R\) и \(I-R\) получается верхнетреугольный блок. Таким образом, весь спектр \(A(u)\) лежит в круге \(|\lambda|<2\), а на \(|\lambda|=2\) резольвента имеет ограниченную норму. Интегральная формула для степени матрицы даёт
\[
                         \|A(u)^{n-1}\|\le C_{M,L_0}\,2^n.     \tag{5}
\]
Оценка действует и для ненормальной \(A(u)\), что важно для открытых сечений у отражений.

Если два выбранных сравнения дают по меньшей мере \(k\) совпадений, число отмеченных \(u\) регулярных сечений не меньше \(\eta n\) при больших \(m\), где \(\eta=(2d-1)/2>0\). В самом деле, без перегруженных сечений два сравнения вместе дали бы не больше \(vn+O_{L_0}(1)\) совпадений. Одно перегруженное сечение добавляет не больше \(v\), а фактическое число совпадений не меньше \(2dm\), \(m=vn+O_{L_0}(1)\).

Из (2), (5) оценка Коши на \(|u|=M, |z_f|=1\) даёт
\[
 \left|\sum_{(x,y)\in E_I}a(x)\overline{a(y)}\right|
 \le C_{M,L_0}\,N^{12L_0}(m+1)^{|I|}\,
                \frac{2^n M^{-\eta n}}{1-1/M}.               \tag{6}
\]
В сумме по степеням \(z_f\) есть не более \((m+1)^{|I|}\) вариантов. Извлечение игнорирует степени \(u<\eta n\), поскольку их вклад в нужные степени \(z_f\) равен нулю. При \(m\ge\lceil\log N\rceil\), \(n\ge m/(2L_0)-2\), правую часть можно сделать \(O(e^{-Am})\) для любого заранее заданного \(A>0\), выбрав \(M\) достаточно большим и затем \(N\) достаточно большим. Все константы остаются равномерными по \(I\).

\newpage
\sect{6. Итог по всем совмещениям и исходному весу}
Точное включение--исключение по \(2m\) сравнениям даёт
\[
 H_{m,N}(d)
   =\sum_{\emptyset\ne I}(-1)^{|I|+1}
                        \sum_{(x,y)\in E_I} a(x)\overline{a(y)}.
\]
Одноэлементные \(I\) дают \(m+O(e^{-c_0 m})\) по разделу 1.
Множества с большим относительным порядком оценены по модулям отдельных пар в разделе 2. Для остальных многомерных пересечений применяем (6) и тот факт, что их число полиномиально по \(m\) при фиксированном \(d\).
Итак, ошибка полного объединения экспоненциально мала. Этот шаг не использует оценку модуля \emph{отдельных} амплитуд в пересечениях малого порядка: там существенна предварительная знаковая сумма.

В частности, для любых фиксированных \(1/2<d_1<d_2\le1\) сумма по парам, у которых максимальная доля совпадений лежит в \([d_1,d_2)\), экспоненциально мала как \(H_{m,N}(d_1)-H_{m,N}(d_2)\). Так контролируется и весь добавленный к прежнему порогу слой без подсчёта его пар по модулям.

По той же схеме объединение, содержащее только обратные сравнения, имеет модуль \(O(e^{-Am})\) при любом \(A>0\). Для одиночных обратных сравнений разрешаем произвольно большой радиус Коши, поскольку все фиксированные \(P V_-^j P\) стремятся к нулю. Для объединения только прямых сравнений сохраняется главный член \(m\).

Циклические пары входят в \(H_{m,N}(d)\) при всяком \(d\le1\). Их точный вклад \(\sum_x d(x)|a(x)|^2\), где \(d(x)\) минимальный период, отличается от
\(m\Tr[(|U_N|^{\circ2})^m]\) не более чем на \(m^2 K^m N^{-m/2}\):
непримитивных слов не больше \(m N^{m/2}\).
С другой стороны,
\(\left|\Tr[(|U_N|^{\circ2})^m]-1\right|\le N(1/100)^m\).
На диапазоне \(m\ge\lceil\log N\rceil\) обе ошибки экспоненциально малы. Это доказывает утверждение и после точного удаления циклических пар.

Для исходного веса \(w\) получаем
\[
 \frac1{N^2}\sum_{\lceil\log N\rceil\le m<2N}w(m/N)H_{m,N}(d)
            \longrightarrow \int_0^2 t w(t)\,dt .
\]
До \(\lceil\log N\rceil\) вклад \(|\Tr U_N^m|^2\) исчезает по третьему входному факту раздела 1, а суммарный вес вычитаемой CUE-рампы равен
\(O((\log N)^2/N^2)\). Отсюда заявленная формула для \(\Delta_{2,N}(b_*)\).
Нужный отрицательный предел остатка равен \(-\delta_*\).

Можно также выбрать \(d=d_N\downarrow1/2\), с неэффективно медленной скоростью. Для последовательности фиксированных \(d_j=1/2+1/(j+2)\) возьмём возрастающие пороги по \(N\), чтобы взвешенная ошибка каждого доказанного результата после его порога была меньше \(1/j\). На промежутках между порогами используем соответствующее \(d_j\). Это даёт ту же формулу для \(\Delta_{2,N}\) с одним меняющимся порогом.

\newpage
\sect{7. Граница полученной оценки и проверки}
\subt{Ровно половина совпадений}
При \(d=1/2\) два сравнения могут набирать ровно по половине без единого тройного совпадения. Например, пусть \(m=2n\) и сравниваются тождество и сдвиг на \(n\). В сечении \((x_j,x_{j+n};y_j,y_{j+n})\) выберем две разные буквы у \(x\). У \(y\) берём либо тот же порядок букв, либо переставленный. При чётном \(n\) используем каждый порядок в \(n/2\) сечениях. Тогда каждое сравнение имеет ровно \(m/2\) совпадений, и в каждом сечении выполнено ровно одно совершенное паросочетание. Поглощение троек здесь не действует. В разделах 8 и 9 этот класс оценён с помощью другой матрицы переноса.

Для типичных множественных пересечений при меньшем пороге требуются суммарные амплитуды переходов между разными парными схемами на последовательных сечениях. Даже контроль пересечений на границе сам по себе не определит полный остаток \(R_{m,N}(d)\).

\subt{Разные пороги у нескольких сравнений}
Для любого фиксированного числа \(k_0\ge2\) различных сравнений в группе ограниченной ширины из раздела 2 доказательство охватывает отдельные пороги \(d_f>0\), если \(\sum_f d_f>1\). В регулярном сечении объединение всех \(k_0\) паросочетаний имеет различные рёбра. Уже \(v+1\) выполненных равенств вынуждают тройку. Каждое перегруженное сечение прибавляет не больше \((k_0-1)v\) совпадений сверх \(v\), поэтому их доля ограничена снизу положительной постоянной. В (3) теперь берётся объединение всех \(k_0\) наборов рёбер, а (4)--(6) остаются верны. Например, знаковая сумма пересечения трёх близких по порядку сравнений с порогом \(d>1/3\) экспоненциально мала. Это отдельная оценка многомерного пересечения; объединение всех сравнений при таких \(d\) здесь ещё не вычислено.

\subt{Проверяемые равенства}
Включённый \texttt{alignment/check.py} проверяет перестановки нитей для прямого и обратного сложения, включая обе чётности временных периодов. Перебор малых матриц сравнивает полный многочлен отметок до и после сложения. Конечные проверки используют точную арифметику корней единицы для полной матрицы Фурье. Отдельный тест сравнивает норму встречи с тройкой со строгой границей \(1/\sqrt N\) для тензора нормированных матриц Фурье. Эти проверки контролируют равенства и направление смежных матриц; оценка для реального \(U_N\) и гейзенберговских времён следует из разделов 1--6.

\subt{Входные результаты в архиве}
Матрицы \(U_N\), тестовый вес и требуемый предел определены в \texttt{research/problem.md}, формулы (2.1)--(2.5), (6.1)--(6.4), (7.11). Оценки всех фиксированных \(j\) в разделе 1 используются из \texttt{contacts/BGS\_Contacts.tex}, разделы 4, 4а и 5. Короткая полоса взята из \texttt{continuation/BGS\_pressure.tex}, формула (1). Прежнее снятие кратности при \(d>2/3\) находится в разделе 6 главы о контактах. Новое доказательство расширяет этот конкретный результат, сохраняя неизменными исходную карту, силу кика и нормировку.

\newpage
\sect{8. Две парные схемы на границе}
Пусть \(m=2n\) и рассмотрим два прямых сравнения с разностью сдвигов \(n\). Суммируем амплитуды пар, давших \(\ge n\) совпадений при каждом сравнении, при условии, что ни в одном сечении из четырёх нитей нет трёх равных букв. В сечении всего два возможных совершенных паросочетания. Так как каждое сечение даёт не более двух равенств, \(\ge2n\) равенств в сумме вынуждает одно из двух полных спариваний на каждом сечении. Буквы \(x\) в сечении различны. Каждая схема используется ровно \(n/2\) раз. При нечётном \(n\) класс пуст.

Положим \(P_{ax}=|(U_N)_{ax}|^2\). На пространстве упорядоченных пар разных букв с размером \(q=N(N-1)\) обозначим через \(D_0\) сжатие \(P\otimes P\). При сохранении парной схемы действует \(D_0\). При переключении действует матрица
\[
   C_{ab,xy}=U_{ax}U_{by}\overline{U_{ay}U_{bx}},
                       \qquad a\ne b,\quad x\ne y.
\]
Пусть \(e_q=q^{-1/2}{\bf1}\). Обе строки и столбцы \(P\) суммируются в 1. При \(a\ne b\) сумма всей строки \(C\) до исключения \(x=y\) равна нулю по ортогональности строк \(U_N\). Аналогично для столбцов. Сумма исключённых членов равна \(\sum_x P_{ax}P_{bx}\le K/N\) для строки и аналогичной сумме для столбца. Для \(D_0\) полная сумма строки и столбца равна 1, а исключённые члены те же. Действительно,
\[
\begin{aligned}
 \langle e_q,D_0e_q\rangle&=1-\theta_N, &
 \langle e_q,Ce_q\rangle&=-\theta_N,\\[-1pt]
 \theta_N&=\frac{N-\sum_{a,x}P_{ax}^2}{N(N-1)}
                      =N^{-1}+O(N^{-2}).
\end{aligned}                                                       \tag{7}
\]
В обеих матрицах нормы блоков между \(e_q\) и его дополнением равны \(O(1/N)\), в том числе при переходе в обратном направлении по оценкам сумм столбцов.

Есть также равномерная оценка \(\|C\|\le\sqrt{K/N}\). До ограничения \(a\ne b,\ x\ne y\) положим \(v_{ab}(x)=\overline{U_{ax}}U_{bx}\). Тогда \(C X\) имеет элементы \(v_{ab}^*Xv_{ab}\) для \(X\) с двумя буквенными индексами. Имеем \(\sum_{a,b}v_{ab}v_{ab}^*=I\) по унитарности и \(\|v_{ab}\|^2\le K/N\). По Коши
\[
 \sum_{a,b}|v_{ab}^*Xv_{ab}|^2\le
        (K/N)\sum_{a,b}\|Xv_{ab}\|^2=(K/N)\|X\|_{\mathrm{HS}}^2 .
\]
Сжатие сохраняет оценку. У \(D_0\) норма на дополнении \(e_q\) не превосходит \(\|P-J_N\|<1/100\). Это следует из разложения \(P\otimes P\) по константам и их дополнениям и последующего сжатия на различные буквы.

\newpage
\sect{9. Коэффициент переключений при половине совпадений}
Введём \(z\), отмечающий вторую схему, и матрицу переноса
\[
       T(z)=\begin{pmatrix}D_0&C\\ zC&zD_0\end{pmatrix}.
\]
На \(|z|=1\) она является сжатием унитарного переноса по четырём нитям, дополненным диагональной унитарной отметкой. В конце сложенного отрезка обе пары нитей переставляются. На внутренних координатах обеих схем эта перестановка есть \(B(a,b)=(b,a)\). На их прямой сумме обозначим её через \(\mathcal B\). Она коммутирует с \(T(z)\) и тождественна на двух копиях \(e_q\). По разложению замкнутого слова искомая сумма равна
\[
                     [z^{n/2}]\Tr\big(\mathcal B\,T(z)^n\big).
\]
Пусть \(T_{00}(z)\) — сжатие на две копии \(e_q\), \(T_{11}(z)\) — на их дополнение, \(E\) — два внедиагональных блока. По разделу 8
\[
                \|T_{11}\|\le \tfrac14,\qquad \|E\|=O(N^{-1})
\]
при больших \(N\), равномерно на \(|z|=1\).
Рассмотрим \(T_s=T_{00}\oplus T_{11}+sE\), \(0\le s\le1\). Это выпуклая комбинация двух сжатий нормы \(\le1\), а её внедиагональные блоки в любой положительной степени имеют норму \(O(N^{-1})\) равномерно по степени. Для блока из малого подпространства в большое раскладываем по последнему переходу: до конца после него остаётся степень \(T_{11}\), и сумма норм не больше \(\|E\|\sum_{j\ge0}4^{-j}\). Для другого направления применяем тот же довод к сопряжённому оператору.

По коммутированию с \(\mathcal B\) производная \(\Tr(\mathcal B T_s^n)\) равна
\(n\Tr(\mathcal B T_s^{n-1}E)\). В следе каждое слагаемое содержит два внедиагональных блока и имеет ранг не больше двух. Поэтому производная есть \(O(n/N^2)\). Интегрируя и оценивая след степени большого блока, получаем
\[
       \Tr(\mathcal B T(z)^n)=\Tr T_{00}(z)^n+
                   O(n/N^2+N^2\,4^{-n}) .
\]
При \(m\asymp N\) остаток \(O(1/N)\) верен и для нужного коэффициента по Коши.

По (7) у \(T_{00}(z)\) переход без смены схемы имеет модуль \(1-\theta_N\), а со сменой — \(\theta_N\le C/N\). Замкнутая бинарная последовательность с обеими схемами содержит \(2j\ge2\) смен. Для коэффициента \(z^k\), \(0<k<n\), число таких последовательностей не больше
\[
       2n\,\binom{k-1}{j-1}\binom{n-k-1}{j-1}.
\]
Выбираем отмеченное начало серии и положительные длины \(j\) серий каждого вида. Суммируя амплитуды по модулям, получаем \(O(n/N^2)\), поскольку при \(n/N\) ограниченном ряд
\(\sum_{j\ge1}(C n/N)^{2j-2}/((j-1)!)^2\) ограничен. Это доказывает оценку \(O(N^{-1})\) на выбранном пересечении. При данном \(m\) у таких пар сравнений \(\le m\) возможных начал, поэтому их вклад после весовой нормировки \(N^{-2}\) стремится к нулю на любой полосе \(m\asymp N\). Остальные пересечения при \(d=1/2\) не оценены этой формулой.

\end{document}
